\chapter{Romax}
\section{Introduction}
    Romax Technology supplies comprehensive simulation software for designing, analyzing, and 
    optimizing drivetrains and rotating machinery across automotive, wind energy, aerospace, and 
    marine sectors. Its tool suite covers durability testing, gearbox and bearing design, structural 
    analysis, and early stage performance rating, enabling rapid, intuitive modelling that improves 
    efficiency, reliability and overall mechanical performance. Romax is recognized as a trusted 
    partner for end-to-end drivetrain development, providing specialized capabilities for electric 
    machine structural design and full powertrain optimization.
    	\begin{figure}[h]
			\centering
			\includegraphics[width=0.6\textwidth]{gearbox/gearbox7}
			\caption{Romax}
			\label{fig:gearbox7}
		\end{figure}
		\FloatBarrier
    
\section{Software Definition}
    The Romax Software Suite is a physics-based Computer-Aided Engineering (CAE) platform built 
    for modeling, simulating, and optimizing rotating machinery and electrified powertrains. 
    Moving away from traditional isolated component design, Romax uses a holistic system 
    architecture. This methodology captures the cumulative, non-linear physical interactions 
    across an entire transmission system under dynamic operational loads, considering the 
    simultaneous effects of component deflections, gear misalignments, shifting contacts, 
    housing flexibility, and non-linear bearing stiffness.

\section{The Romax Hybrid System-Simulation Approach}
    To accurately analyze these complex interactions, Romax software uses a hybrid modeling 
    approach that integrates multiple computational methods into a single, coupled simulation, 
    balancing calculation speed with high prediction accuracy:
        \begin{itemize}
			\item \textbf{Finite Element (FE) Method:} Condenses large structures, such as flexible gear housings or planet carriers, into reduced substructures to capture structural stiffness and mass distribution efficiently.
			\item \textbf{Analytical Formulations:} Computes shaft bending lines and gear meshing mechanics using closed-form analytical equations and advanced gear/bearing contact models.
			\item \textbf{Empirical Formulations:} Estimates bearing clearances and operational fatigue lives according to standard codes, such as international rolling bearing rating life methods.
		\end{itemize}

\section{Features and Key Benefits}
    Romax offers a comprehensive CAE solution for drivetrain and powertrain engineering, 
    delivering cutting-edge simulation and analysis capabilities that empower engineers to design, 
    optimize, and validate complex systems with precision and speed.

    The product delivers a holistic system view, capturing intricate component interactions and 
    enabling design of the entire system from physics based, parametric definitions for every part; 
    it relies on proven, proprietary technology that intelligently blends analytical, empirical, and 
    finite element approaches, coupled with world-class gear and bearing contact models. By allowing 
    models to be built and evaluated in minutes, it delivers speed and simplicity, providing rapid, 
    actionable results that let engineers run “what if” scenarios and gain instant insight. 
    It encourages CAE-driven design by merging easily editable models with repeatable, multi-fidelity 
    analyses and multi-attribute optimisation, empowering better engineering decisions and reducing 
    innovation risk. The flexible product suite streamlines processes with collaborative workflows, 
    optimization, and automation that boost productivity. Integrated into the broader Hexagon 
    ecosystem, the tool breaks organizational barriers through seamless connections between CAE, 
    CAD, and FEA, specifically enhancing powertrain design and beyond.

\section{Integrated Applications}
    The platform architecture utilizes six distinct, interoperable modules engineered to support 
    collaborative engineering workflows across cross-functional powertrain development teams:
        \begin{itemize}
			\item \textbf{Romax Enduro:} The primary structural durability interface used to perform fatigue evaluations, component structural safety factor checks, and system deflection analysis.
			\item \textbf{Romax Spectrum:} A highly specialized vibration and acoustics module designed to model dynamic transmission error (TE), predict structural noise paths, and map housing acoustic radiation. Also, it facilitates the simulation of noise and vibration (NVH) in electro-mechanical powertrains.
			\item \textbf{Romax Spin:} An advanced internal element modeling workspace allowing bearing specialists to map roller-edge stress concentrations, local contact footings, and internal contact geometry.
			\item \textbf{Romax Concept:} A streamlined workspace optimized for clean-sheet structural layout definition and spatial envelope exploration. Also, it can help investigate EV drivetrain concepts in their early stages.
			\item \textbf{Romax Energy:} A dedicated efficiency and fluid-loss solver that isolates and quantifies gear mesh friction, rolling element resistance, and hydrodynamic oil churning losses.
			\item \textbf{Romax Evolve:} An electro-mechanical simulation module that couples stator and rotor structural dynamics with multi-phase electromagnetic forces in unified electric vehicle powertrains. Also, it allows for structural design and analysis of electrical machines.
		\end{itemize}

        \begin{figure}[h]
			\centering
			\includegraphics[width=0.9\textwidth]{gearbox/gearbox8}
			\caption{Romax}
			\label{fig:gearbox8}
		\end{figure}
		\FloatBarrier

\section{Automated System Sizing Interfaces}
    To transition a preliminary generative design into a validated powertrain, the structural 
    boundaries established by the closed-form weight and sizing equations are translated directly 
    into the physics-based environment of the Romax Concept. Rather than operating as an independent, 
    isolated component model, the software uses a multi-layered parametric data pipeline.

    \subsection{Romax Primary System Inputs}
        The analytical model requires a rigorous data matrix defining the full macro-geometric, 
        kinematic, and operational layout of the transmission before it can execute any non-linear 
        structural sweeps:
            \begin{itemize}
                \item \textbf{Operational Conditions and Duty Cycles:} The software requires explicit multi-step loading profiles defining nominal input power, engine input velocity, and target main rotor mast velocity.
                \item \textbf{Gear Macro-Geometry:} The system architecture requires the exact parameters derived from the optimization loop: standard metric module $m_{sta}$, pressure angle $\alpha_{n}$, helix angle $\beta$, and specialized spiral bevel tooth indices. Sizing profiles must contain explicit teeth counts—the sun teeth $z_{s}$, planet teeth $z_{p}$, and ring teeth $z_{r}$—along with the precise face widths $F_{b}$ and $F_{p}$.
                \item \textbf{Planetary Gearbox Boundaries:} The spatial inputs include the precise number of active planet gears and the geometric layout positions of the planet pins relative to the vertical centre line.
                \item \textbf{Shaft Geometry:} Shaft inputs discard one-dimensional line assumptions, requiring complete 2D/3D profile coordinate matrices. The pipeline defines explicit shaft length segments along with outer shaft step diameters, internal hollowing parameters for weight reduction, and the physical dimensions of the planet pins.
                \item \textbf{Rolling Element Bearing Selection:} Sizing requires either direct lookup selection from integrated commercial manufacturer databases (e.g., SKF, NSK, Timken catalog lines) or explicit definition of internal rolling geometry. This includes roller/ball counts, raceway diameters, pitch circles, internal clearances, and the specific structural preload vectors acting on the input and output main line supports.
                \item \textbf{Material Fatigue and Lubricant Limits:} Sizing require exact core and case material properties, including densities, Young's Modulus, Poisson's ratio, and specific allowable fatigue limits. Fluid mechanics constraints require the lubricant type, dynamic viscosity indices, operating temperatures, and oil-jet delivery pressures.
            \end{itemize}

    \subsection{Romax Multi-Disciplinary Outputs}
        Once the parametric input is compiled, the solvers execute a whole-system structural study. 
        The software returns four primary engineering evaluation modules:
            \begin{itemize}
                \item \textbf{Stress Analysis and Component Safety Factors:} The system outputs exact, code-validated safety margins. It logs individual safety factors for tooth root bending, surface contact pressure, and shaft torsional shear fatigue.
                \item \textbf{Whole-System Mass Bookkeeping and Spatial Size:} Romax treats all gears, shafts, and casing components as true 3D solid geometries. It integrates the exact volume of the thin-walled housing and the webbed gear wheels to output the Total Assembled Weight.
            \end{itemize}

\section{Romax in AI/ML Based Transmission Design}
    The paper shows how Romax's physics based; full system gearbox simulation platform can be combined 
    with ODYSSEE's AI/ML tools to create a practical, high-performance optimization loop. It uses AI 
    to accelerate and guide the traditional physics-based simulation workflow.

    It isn't exactly what the project will do, but it gives an idea of how Romax will be used and how 
    AI improves traditional calculations. The concept is like the GAIT project: instead of calculating 
    gearbox weight and size one by one with many formulas, AI assists and replaces traditional hand 
    calculations. With Romax, we can also verify the results.

    It outlines a shift in transmission design driven by sustainability, electrification, and shorter 
    development cycles. While conventional CAE is mature, AI/ML can accelerate the simulation of 
    lifecycle, democratize advanced engineering, and reduce risk. The authors fuse physics-based 
    CAE with AI/ML, moving from simulation test validation to Design-of-Experiments fed machine 
    learning models verified by simulation and experiment. Two case studies illustrate this: 
    a gear-micro-geometry optimization that reduces transmission error, and a larger study that 
    iteratively optimizes numerous design parameters against broad performance criteria. The method 
    applies beyond automotive—into aerospace and other sectors—enabling full design space exploration, 
    stronger multiphysics optimization, and ensuring gearboxes match simulated behavior. The ultimate 
    aim is faster, higher quality powertrain development with less reliance on specialist expertise.

        \begin{itemize}
            \item \textbf{Romax}\par
                Parametric, hybrid FEM or empirical or analytical simulation of entire gearboxes 
                (gears, shafts, bearings, CVT, etc.). Fast (<1s per run) and highly accurate for gearbox 
                mechanics, durability, NVH, efficiency, weight plus packaging.
            \item \textbf{ODYSSEE}\par
                AI / ML platform that builds Reduced-Order Models (ROMs) from a handful of selected 
                simulation runs and then performs design-of-experiments (DoE), sensitivity analysis, 
                and optimisation on that ROM.
            \item \textbf{Romax to ODYSSEE to Romax}\par
                \begin{enumerate}
                    \item Define the full system model in Romax (parametric).
                    \item Use Romax batch-run to generate a small training data set.
                    \item Train a ROM in ODYSSEE.
                    \item Run millions of parametric evaluations on the ROM (but actually far fewer than would be needed in Romax).
                    \item Identify the best parameter combinations.
                    \item Re-run the selected candidates in full Romax to validate and finalize the design.
                \end{enumerate}
        \end{itemize}